% Transcription — fonds Alexandre Grothendieck, shelfmark 70, batch 3 (pages 41-60).
% Established from the facsimile archives/batches/70/batch-03.pdf (cover sheet
% excluded; batch page 1 = archivists' page 41) and the 700-dpi tiles in
% archives/tiles/70/p41 ... p60, per the skill transcribe-grothendieck.
% Folder 70 is « Réalisations géométriques de structures combinatoires
% (n-polyèdres, n-hyperpolyèdres [2-polyèdres réguliers]…) (Vieilles
% rédactions) : notes manuscrites (s.d.) », 125 pages; the bracketed part of
% the title is the archivists'. Read in parallel with the other batches of the
% folder and without consulting them: notation fixed from these pages alone.
%
% Content: pp. 41-47 continue, from before this batch, a draft on n-cartes
% (Coxeter-type group G_n = <sigma_0..sigma_n | sigma_i^2, (sigma_i sigma_j)^2
% for |i-j|>=2>) and their geometric realisations by maximal affine flags;
% p. 41 opens mid-argument and p. 47 breaks off mid-sentence. pp. 49 and 50 are
% his covers (« 2 – Polyèdres réguliers. (Vieille rédaction) », « (18)
% 2-Polyèdres repérés en coordonnées universelles (vieille rédaction) »),
% rendered as headings. pp. 51-56: geometric and regular realisations of a
% combinatorial 2-polyhedron X = (S, A, F, R, R', R''). p. 58: a separate sheet
% on groups acting on a groupoid / graph. p. 59: figures only.
%
% Pages skipped: 48 (blank), 57 (a draft letter in his hand, cancelled by two
% diagonals — not mathematics), 60 (blank but for faint offset of another
% sheet's text). Pages 49-50 carry only his cover titles and get no \page.
%
% Notation: his underlined letters for the sets of the structure and for
% elements of the automorphism group are set in roman (\mathrm{S}, \mathrm{g},
% \mathrm{G}); plain italics for their images in E.
%
% Pass: Opus 5.5 (claude-opus-5-5), 2026-09-29 — first pass, unchecked against the pages by a human.

\documentclass[11pt,a4paper]{article}
\input{../preamble/grothendieck}

\folder{70}
\batch{3}
\pages{41}{60}
\dating{[à partir de 1976]}
\watermark{Édition de démonstration}
\foldertitle{Réalisations géométriques de structures combinatoires (n-polyèdres, n-hyperpolyèdres [2-polyèdres réguliers]…) (Vieilles rédactions) : notes manuscrites (s.d.).}

\begin{document}

\page{41}
\note{la page s'ouvre au milieu d'un argument commencé avant ce lot : $E$, les $\sigma_i$, les groupes $G^{-}(i)$, $G^{+}(j)$ et les $D_{i,j}$ sont définis plus haut.}
\[
E(i,j) = E/(\sigma_0,\dots,\sigma_{i-1};\ \sigma_{j+1},\dots,\sigma_n) = E/G^{-}(i)\times G^{+}(j) = D_{i,i+1,\dots,j}
\]
\note{le dernier membre, $D_{i,i+1,\dots,j}$, est ajouté par lui au-dessus de la ligne.}
et l'application
\[
E(i,j) \longrightarrow D_{i,j} .
\]
\struck{\ill{}} \add{$0\leq i\leq i+2\leq j\leq n+1$}
Le groupe engendré par $\sigma_{i+1},\dots,\sigma_{j-1}$ opère sur $E(i,j)$, avec comme quotient $D_{i,j}$. \ill{}, donc il opère \uncertain{cartographiquement}\note{souligné par lui}, et on trouve une $(j-i-2)$-\ill{}, dont les composantes connexes correspondent aux $\Pi(d_{ij})$, avec $d_{ij}\in D_{ij}$ ; les \uncertain{facettes de dim} $h$ ($0\leq h\leq j-i-2$) correspondent aux drapeaux de type $(i,\underline{i+h+1},j)$, $i+1\leq\ \leq j-1$, au-dessus des drapeaux $d_{ij}$.

Cas particuliers : $j=i+2$ \uncertain{un pas idiot}\note{lecture très incertaine} ;

$j=i+3$ : on trouve un \uncertain{polygone}, d'un invariant $p(d_{ij})\in\overline{\mathbf{N}^{*}}$ (où $1\leq p(d_{ij})\leq+\infty$).

\struck{Des \ill{} au} \struck{Introduisons} Formellement les $\Pi^{*}(f_i)$ correspond\supplied{ent} \struck{\ill{}} \struck{\ill{}} $\Pi(d_{-1,i})$, $d_{-1,i}$ défini par $f_i$ \struck{(\ill{})}
\struck{(-1) \ill{}} et $\Pi^{*}(f_i')$ \ill{} à $\Pi(d_{i,n+1})$, où $d_{i,n+1}$ est défini par $f_i$. Les invariants \uncertain{essentiels}

\page{42}
essentiels sont
\begin{enumerate}
\item[a)] Les $\nu(\Pi^{-}_{f_2})$ --- disons $p_0(\Pi)$ si $\Pi$ régulier ;
\item[b)] Les $\nu(\Pi(d_{0,3}))$, $\nu(\Pi(d_{1,4}))$, \ldots\ ($\nu(\Pi(d_{n-3,n}))$) --- notés $p_1(\Pi)$, $p_2(\Pi)$, \ldots, $p_{n-2}(\Pi)$ si $\Pi$ régulier ;
\item[c)] Les $\nu(\Pi^{\rightarrow}_{f_{n-2}})$, \struck{\ill{}} notés $p_{n-1}(\Pi)$ si $\Pi$ régulier.
\end{enumerate}

Supposons $\Pi$ régulier, je dis que
\[
\begin{cases}
p_0 = \text{ordre de } \sigma_0\sigma_1\\
p_1 = \text{ordre de } \sigma_1\sigma_2\\
\quad\dots\\
p_{n-1} = \text{ordre de } \sigma_{n-1}\sigma_n
\end{cases}
\]
\marginal{vrai tout au moins dans le cas des géométries d'incidence (\og polyèdres\fg\ldots)}

En effet, \struck{\ill{}} sur $\Pi$ a régularité, $\nu(\Pi^{\rightarrow}_{f_2})$ est l'ordre de $\sigma_0\sigma_1$ opérant sur
\[
\bigl(E/G^{+}_{n}(2)\bigr)_{f_2} \quad\text{ou}\quad E^{-}(2)_{f_2}
\]
\struck{ou $\Pi(d_{i,i+3})$ est l'ordre}\note{dans le passage biffé, le premier indice de $d$ est illisible ; $i$ est une lecture par défaut.}
\[
(\sigma_0\sigma_1)^{p}x = g_x x \qquad g_x\in G^{+}_{n}(2)
\]
\drawing{petit schéma d'indices : $i-1$ ; $i$, $i+1\leq\ \leq j-1$, $j+1$ ; $j\geq i+2$ ; puis un $i$ isolé.}
\note{le reste de la page est vide.}

\page{43}
$\Pi$ une $n$-carte \add{\uncertain{ou symb.}} (\ill{} \ill{} $R=\mathrm{Rep}(\Pi)$, \ill{}\note{le coin supérieur droit du feuillet manque ; la fin de la ligne est perdue.}
$\mathfrak{G}_n$-ensemble (où
\[
\mathfrak{G}_n = \bigl\{\sigma_0,\dots,\sigma_n \bigm| \sigma_0^2=\dots=\sigma_n^2=1,\ \sigma_i\sigma_j=\sigma_j\sigma_i \text{ si } j\geq i+2\bigr\}\ )
\]
Une \emph{représentation géométrique} $\varphi$ de $\Pi$ dans \struck{\ill{}} un espace affine $E$ (sur un corps \uncertain{ordonné}) \struck{\ill{}} est la donnée d'une application
\[
(1)\qquad R \xrightarrow{\ \varphi_R\ } \mathrm{Drap\,max}(E) = \bigl\{(E_0\subset E_1\subset\dots\subset E_i\subset\dots\subset E_n\subset E)\bigr\}
\]
\marginal{sous-espace affine de $E$ de rg $i$ (renvoi sur $E_i$)}
telle que pour $r\in R$, $0\leq i\leq n$, on ait
\[
(2)\qquad
\begin{cases}
\pi_j(\varphi_R(\sigma_i r)) = \pi_j\varphi_R(r) \quad\text{si}\\
0\leq i,j\leq n,\ i\neq j
\end{cases}
\]
où $\pi_j(d)$ désigne la composante de dim $j$ d'un drapeau $d$. Je veux \ill{} \add{dire par} passage au quotient, $\pi_j\varphi_R$ se factorise en
\[
(3)\qquad D_j(\Pi) \xrightarrow{\ \varphi_j\ } \mathrm{Grass\,aff}_j(E) \qquad (0\leq j\leq n)
\]

\page{44}
\note{toute la page est barrée de deux grandes diagonales ; elle se lit néanmoins, et c'est un premier essai de la définition de la p.~43.}
Considérons une réalisation géom\supplied{étrique} $\varphi$ de $\Pi$ \add{\uncertain{dans un espace affine $E$ de dim $n+1$}} \marginal{(\uncertain{régulière} \uncertain{épingl.}) ?}
\uncertain{(affine disons)} sur un anneau $k$, donnée par

\begin{tikzcd}
G=\mathrm{Aut}(\Pi) \arrow[r, "\varphi_G"] & \mathrm{Aff}(E) \\
R=\mathrm{Rep}(\Pi) \arrow[r, "\varphi_R"] \arrow[d] & \mathrm{Drap\,max}(E) \arrow[d, "\pi_i"] \\
D_i(\Pi) \arrow[r, "\varphi_{D_i}"] & \mathrm{Grass}_{i}(E)
\end{tikzcd}

\note{sous $R$ : « $G$-torseur », sous la flèche $\varphi_R$ : « $G$-hom\supplied{omorphisme} » ; l'indice de $\mathrm{Grass}_{i}$ est surchargé et douteux.}
\[
\begin{cases}
\varphi_R(r) = (E_0\subset E_1\subset\dots\subset E_n\subset E)\\
\sigma_i E_j = E_j \quad\text{si } i\neq j \qquad \dots
\end{cases}
\]

\page{45}
Lorsque la $n$-carte $\Pi$ est un polyèdre (une géométrie d'incidence), dans la donnée de $\varphi_R$ satisfaisant (2) équivaut : à la donnée des $\varphi_j$ ($0\leq j\leq n$), soumis\supplied{es} à la condition d'être compatibles aux relations d'incidence. Dans le cas \add{encore} général, on peut encore dire que (3) se généralise en des
\[
(4)\qquad D_{\{j_0,j_1,\dots,j_p\}}(\Pi) \xrightarrow{\ \varphi_{\{j_0,\dots,j_p\}}\ } \mathrm{Drap\,aff}_{\{j_0,j_1,\dots,j_p\}}(\Pi)
\]
\note{le but de la flèche est écrit $\mathrm{Drap\,aff}_{\{j_0,\dots,j_p\}}(\Pi)$ ; on attendrait $(E)$.}
pour toute partie $\{j_0<j_1<\dots<j_p\}$ de $\Delta_n=\{0,1,\dots,n\}$, ces applications étant compatibles aux applications de restriction \uncertain{évidentes} par les \uncertain{inclusions} de parties de $\Delta_n$. La donnée de $\varphi_R$ satisfaisant (2) équivaut à la donnée des applications (4) satisfaisant la condition de fonctorialité précédente \ill{} les inclusions \add{\uncertain{les $\varphi_R$}} (\uncertain{comme} \ill{} on peut donc \uncertain{les} \ill{} des \uncertain{hom.}\ de géométries de drapeaux \add{l'une et l'autre} \uncertain{basées} sur $\Delta_n$, de $D(\Pi)$ dans $\mathrm{Drap\,aff}(E)$).

\page{46}
\note{la première moitié de la page (jusqu'à « Fixons une telle composante connexe ») est barrée de deux diagonales ; elle reste lisible.}
\struck{Ceci dit, considérons la $\nu=(j-i-2)$-carte \add{$\Pi_{i,j}$} définie par le $\mathfrak{G}$-ensemble \uncertain{quotient} ($0\leq i\leq i+2\leq j\leq n$)}
\[
R_{ij} = R/(\sigma_0,\dots,\sigma_{i-1},\sigma_{j+1},\dots,\sigma_n)
\]
\struck{(NB $R_{0j}=R/(\sigma_{j+1},\dots,\sigma_n)$, $R_{i,n}=R/(\sigma_0,\dots,\sigma_{i-1})$) avec les opérations induites par le groupe \ill{} quotient par}
\[
\begin{matrix}
\sigma_{i+1}, & \dots, & \sigma_{j-1}\\
\shortparallel & & \shortparallel\\
\sigma'_0, & \dots, & \sigma'_{j-i-2=\nu}
\end{matrix}
\]
\struck{qui en fait une $\nu$-carte, dont l'ensemble des composantes connexes s'identifie à $D_{ij}(\Pi)$. Fixons une telle composante connexe,} i.e.\ un

Ceci dit, soit $f_i\in D_i(\Pi)$ \add{($1\leq i\leq n$)}, $E_i=\varphi_i(f_i)$, et considérons la \add{$(i-1)$-carte induite} $\partial f_i$ dont l'ensemble des repères est $D_{01\dots i}(f_i)$
\[
= \text{image réciproque de } \{f_i\} \text{ par } D_{0,1,\dots,i}(\Pi)\to D_i(\Pi).
\]
\add{un \emph{\uncertain{bouquet}} de $\varphi$} les \uncertain{données} \uncertain{définissent} \uncertain{une} réalisation géom\supplied{étrique} canonique (dite \og induite\fg) de $\partial f_i$ dans $E_i$.

Légèrement moins \uncertain{évident} est la définition de la représentation induite de la $(n-i-1)$-carte $\delta f_i$ dans l'un des

\page{47}
\ill{}\note{mot de raccord incertain, la phrase continue la p.~46} dont l'ensemble des repères est $D_{i,i+1,\dots,n}(f_i)=$ image inverse de $\{f_i\}$ par $D_{i,i+1,\dots,n}\to D_i$. En les \uncertain{sommets} sont les \struck{\ill{}} $(i,i+1)$-drapeaux de $\Pi$ dont la $i$-composante est $f_i$, et ils s'envoient sur les drapeaux de $E$ de type $(i,i+1)$ dont la $i$-composante est $E_i$, \struck{\ill{}} i.e.\ des couples $(E_i\subset E_{i+1})$, $E_{i+1}$ de rg $i+1$. Mais si $V_i$ est l'espace des translations de $E_i$, les $E_{i+1}$ correspondent : aux droites de $\overline{E}=E/V_i$
\note{la page s'arrête ici ; le reste est blanc.}

\section*{2 -- Polyèdres réguliers. (Vieille rédaction)}
\note{titre de sa main, en haut à droite d'une feuille de chemise autrement vierge (p.~49) ; la p.~48 est vierge.}

\subsection*{2-Polyèdres repérés en coordonnées universelles (vieille rédaction)}
\note{titre de sa main, au crayon, sur une seconde chemise (p.~50), précédé de « (18) ».}

\page{51}
\note{il souligne les lettres désignant les ensembles de la structure ($\mathrm{S}$, $\mathrm{A}$, $\mathrm{F}$, $\mathrm{R}$, \ldots) et les éléments du repère ; on les rend ici en romain droit.}
$X=(\mathrm{S},\mathrm{A},\mathrm{F},\mathrm{R},\mathrm{R}',\mathrm{R}'')$ polyèdre combinatoire \uncertain{(\ill{})}, $\mathrm{R}$ ens\supplied{emble} de ses repères.

$E$ espace affine (proj.) sur $k$ (\uncertain{corps}).

\emph{Réalisation géom\supplied{étrique}} de $X$ :
\[
\begin{cases}
\varphi_S : \mathrm{S}\to E\\
\varphi_A : \mathrm{A}\to \mathrm{Dr}(E)\\
\varphi_F : \mathrm{F}\to \mathrm{Pl}(E)
\end{cases}
\]
1°) compatible aux $\mathrm{R}'$, $\mathrm{R}''$, donc $\mathrm{R}$. On suppose de plus que

2°) a) Si $a\in\mathrm{A}$, $s,s'\in\mathrm{S}$ les deux sommets incidents, $f,f'$ les deux faces incidentes, on a $\varphi_S(s)\neq\varphi_S(s')$, $\varphi_F(f)\neq\varphi_F(f')$.

b) Si $s<f$ avec $s\in\mathrm{S}$, $f\in\mathrm{F}$, et si $a,a'$ sont les deux arêtes entre $s$ et $f$, on a $\varphi_A(a)\neq\varphi_A(a')$.

$\Longrightarrow$ $\varphi_S$ détermine $\varphi_A$, $\varphi_F$ / dualement $\varphi_F$ détermine $\varphi_A$, $\varphi_S$ /
\add{$\varphi_A$ détermine $\varphi_S$ dualement $\varphi_F$}

\struck{Si $X$ est régulier, on appelle réalisation géom\supplied{étrique}} \emph{régulière}.

\emph{Prop.} Soit $r_0=(s_0,a_0,f_0)$ un repère de $X$. Soit :

$s_1$ le deuxième sommet de $a_0$, $s_2$ le deuxième sommet de la 1\supplied{ère} arête \add{$a_1$} entre $s_0$ et $f_0$, $s_3$ le deuxième sommet de la 2\supplied{ème} arête entre $s_0$ et $f_1$, la 2\supplied{ème} face passant par \uncertain{$a_0$}.

Considérons les images $s_0,s_1,s_2,s_3$ dans $E$. Alors ces pts sont affinement indépendants

\page{52}
\emph{Dém.} $s_1\neq s_0$ (par a) 1°) \add{et $s_2\neq s_0$}\,; $a_1\neq a_0$ par b), donc $s_2\notin a_0$, donc $s_0,s_1,s_2$ aff\supplied{inement} indép\supplied{endants}, $f_0$ est le plan qu'ils engendrent. De même [remplaçant $r_0=(s_0,a_0,f_0)$ par $r_1=(s_0,a_0,f_1)$] on voit que $f_1$ est le plan engendré par $s_0,s_1,s_3$. Or $f_1\neq f_0$ (\struck{\ill{}} par a) 2°), donc $s_3\notin f_0$, donc $s_0,s_1,s_2,s_3$ aff\supplied{inement} indép\supplied{endants}.

\emph{Cor.}
\[
\left.\begin{aligned}
a_0&=\mathrm{dr}(s_0,s_1)\\
f_0&=\mathrm{pl}(s_0,s_1,s_2)
\end{aligned}\right\}\Longrightarrow \varphi_A,\varphi_F \text{ connus quand on connaît } \varphi_S
\]
dualement
\[
\left.\begin{aligned}
a_0&=f_0\cap f_1\\
s_0&=f_0\cap f_1\cap f_2
\end{aligned}\right\}\Longrightarrow \varphi_S,\varphi_A \text{ connus quand on connaît } \varphi_F
\]
où $f_2$ est la deuxième face passant par $a_1$, et enfin
\[
\left.\begin{aligned}
s_0&=a_0\cap a_1\\
f_0&=\mathrm{pl}(a_0,a_1)
\end{aligned}\right\}\Longrightarrow \varphi_S,\varphi_F \text{ connus quand on connaît } \varphi_A
\]
\drawing{figure : le sommet $s_0$ d'où partent les arêtes $a_0$ (vers $s_1$), $a_1$ (vers $s_2$), $a_2$ (vers $s_3$, vers le haut) et une quatrième arête vers le bas, aboutissant à un point marqué $s_3$ (surchargé) ; les faces $f_0$ (entre $a_0$ et $a_1$, hachurée), $f_1$ (entre $a_2$ et $a_0$, hachurée), $f_2$ (sous $a_1$) et $f_3$ au-delà de $s_1$ ; une arête $a_2^0$ issue de $s_1$.}

\emph{Conclusion} $\varphi_S(\mathrm{S})$ engendre aff\supplied{inement} $E$.

\emph{Réalisation géom\supplied{étrique}} \emph{régulière} (où $X$ régulier) :
$\forall \mathrm{g}\in \mathrm{G}\overset{\mathrm{df}}{=}\mathrm{Aut}(X)$, $\exists\, g\in\mathrm{Aut}_{\mathrm{aff\ ou\ proj}}(E)$
\note{la phrase se poursuit en p.~53.}

\page{53}
tel que \struck{\ill{}} l'on ait
\[
\begin{cases}
g\circ\varphi_S = \varphi_S\circ\mathrm{g}\\
g\circ\varphi_A = \varphi_A\circ\mathrm{g}\\
g\circ\varphi_F = \varphi_F\circ\mathrm{g}
\end{cases}
\]
\note{au-dessus de chaque $\mathrm{g}$ souligné du membre de droite, un petit signe (accent ou flèche) que je ne sais pas interpréter.}

\emph{NB} l'une quelc\supplied{onque} de ces trois relations implique les deux autres. Dans le cas affine, $g$ défini en fonction unique de $\mathrm{g}$, et $\mathrm{g}\mapsto g$ est un homomorphisme $\mathrm{G}\to$ \uncertain{opérations} sur $E$. C'est encore vrai dans le cas projectif, pourvu que
$\exists\, s\in\mathrm{S}$ avec $\varphi_S(s)=s$ \uncertain{non coplanaire avec} trois des quatre pts $s_0,s_1,s_2,s_3$ (condition
\marginal{à vérifier}
exception\supplied{nelle} : \struck{tétraèdre \ill{}} \struck{\ill{}} $\mathrm{S}=\{s_0,s_1,s_2,s_3\}$\,\ldots) Si \ill{} nous éliminons ce cas exceptionnel, il y a un \emph{\uncertain{donc}} $\mathrm{g}\mapsto g$ \uncertain{faisant} opérer $\mathrm{G}$ sur $E$\,\ldots.

Donnons-nous supposons $X$ régulier, on choisit repère $r_0=(s_0,a_0,f_0)$, d'où générateurs $\sigma_0,\sigma_1,\sigma_2$ de $\mathrm{G}=\mathrm{Aut}(X)$, définis par
\note{la fin de la page est occupée par un mot surchargé sous « générateurs », illisible ; la définition des $\sigma_i$ se poursuit en p.~54.}

\page{54}
\[
\begin{cases}
\sigma_0(s_0)\overset{\mathrm{df}}{=}s_1, & \sigma_0(a_0)=a_0,\quad \sigma_0(f_0)=f_0\\
\sigma_1(s_0)=s_0 & \sigma_1(a_0)\overset{\mathrm{df}}{=}a_1,\quad \sigma_1(a_0)=a_0\\
\sigma_2(s_0)=s_0 & \sigma_2(a_0)=a_0,\quad \sigma_2(f_0)\overset{\mathrm{df}}{=}f_1
\end{cases}
\]
\note{deuxième ligne : il écrit bien $\sigma_1(a_0)=a_0$ en troisième position ; on attendrait $\sigma_1(f_0)=f_0$.}

La donnée de $\varphi$ régulière équivaut : à la donnée de $\varphi_G:\mathrm{G}\to\mathrm{Aut}(E)$ et $\varphi_R(r_0)=(s_0,a_0,f_0)\in\mathrm{Drap}_{012}(E)$, satisfaisant certaines conditions. La donnée de $\varphi_G$ : équivaut en termes \uncertain{équivaut} : à la donnée de $\sigma_0,\sigma_1,\sigma_2\in\mathrm{Aut}(E)$ satisfaisant certaines conditions, qu'on explicitera.

Notons qu'indépendamment de la condition de régularité et de l'épinglage \uncertain{donné} de $r_0\in\mathrm{R}$, la donnée de $\varphi$ implique
\[
\varphi_R : \mathrm{R}\longrightarrow \mathrm{Drap}_{012}(E)
\]
qui, dans le cas \emph{régulier}, devient une $\mathrm{G}$-application, et qui permet de récupérer $\varphi_S,\varphi_A,\varphi_F$ par passage aux quotients. Dans le cas \struck{\ill{}} régulier, comme $\mathrm{R}$ est un $\mathrm{G}$-torseur, la donnée de $\varphi_R$ équivaut : à la donnée de $\varphi_R(r_0)\in r_0=(s_0,a_0,f_0)$).
\note{« $\varphi_R(r_0)\in r_0=(s_0,a_0,f_0)$ » : ainsi sur la page ; le « $\in$ » est peut-être un « $=$ » surchargé.}

\page{55}
Pour que la \struck{\ill{}} $\mathrm{G}$-application \uncertain{\ill{}} $\pi_E^{(i)}\circ\varphi_R$ passe au quotient \add{par} \ill{} \add{$\mathrm{Fac}_i(X)$ ($=\mathrm{S},\mathrm{A}$ ou $\mathrm{F}$), \uncertain{soit} $\mathrm{G}$-\uncertain{équivariante}} i.e.\ qu'on ait

\begin{tikzcd}
\mathrm{R} \arrow[r, "\varphi_R"] \arrow[d, "\pi_X^{(i)}"'] & \mathrm{Drap}_{012}(E) \arrow[d, "\pi_E^{(i)}"] \\
\mathrm{Fac}_i(X) \arrow[r] & \mathrm{Fac}_i(E)
\end{tikzcd}

$i=0,1,2$

\struck{($s,a,f\Rightarrow$ l'image)} le stabilisateur dans $\mathrm{G}$ de $f(i)_0$ contienne celui de $f(i)_0$, ce qui s'écrit
\note{les deux « $f(i)_0$ » sont ainsi écrits, l'un souligné (dans $X$), l'autre non (dans $E$).}
\[
\begin{matrix}
 & \sigma_0(a_0)=a_0 & \sigma_0(f_0)=f_0\\
\sigma_1(s_0)=s_0 & & \sigma_1(f_0)=f_0\\
\sigma_2(s_0)=s_0 & \sigma_2(a_0)=a_0 &
\end{matrix}
\]
\note{au-dessus du tableau, un mot biffé ; le reste de la page est blanc.}

\page{56}
\add{\uncertain{\ill{}},} \emph{À vérifier} :
\uncertain{Cartes homogènes} de type $(p,q)$

a) \add{\uncertain{Rep.}} est géom\supplied{étrie} d'incidence ssi $(p,q)=(1,2)$ ou $(2,1)$ (seuls cas où \uncertain{\ill{}} pour $(p,q)\geq(2,2)$)
\drawing{deux petites figures : un segment à deux sommets marqués ; une courbe fermée portant un seul sommet.}

b) est polyèdre combinatoire ssi $p=2$ ou $q=2$

\struck{Alors}

[C'est une géom\supplied{étrie} d'incidence ssi $(p,q)\geq(2,2)$ i.e.\ $p\neq1$, $q\neq1$ i.e.\ $(p,q)\neq(1,2),(2,1)$]

[C'est un polyèdre combinatoire ssi $(p,q)\geq(3,3)$ i.e.\ ssi $p\neq2$, $q\neq2$, i.e.\ ssi $(p,q)\neq(2,q)$ ($q\geq1$) et $(p,q)\neq(p,2)$ ($p\geq1$).]
\note{les deux énoncés entre crochets corrigent visiblement a) et b), écrits d'abord : a) et b) disent l'inverse (« ssi $(p,q)=(1,2)$ ou $(2,1)$ », « ssi $p=2$ ou $q=2$ ») ; il n'a pas biffé a) et b).}

\page{58}
\note{feuillet isolé, d'une autre encre, sans lien direct visible avec les pp.~51--56 : groupes d'automorphismes d'un groupoïde et d'un graphe.}
$\mathcal{C}$ \uncertain{groupoïde}, $G$ \add{$\subset\mathrm{Aut}(\mathcal{C})$} transitif sur $\mathrm{Ob}\,\mathcal{C}$, $e\in\mathcal{C}_0$, $H=G_e$, $\mathcal{C}_0\simeq G/H$.

$\Gamma$ normalisateur de $G$

$\Gamma^0$ centralisateur de $\mathcal{C}$ \note{sic ; on attendrait « de $G$ ».}
\[
1\to\Gamma^0\to\Gamma\to\mathrm{Aut}(G)
\]
$e\in\mathrm{Ob}\,\mathcal{C}=\mathcal{C}_0$

$\gamma$ \add{$\mathcal{C}_0$} connu quand on connaît $\gamma(e)$


\begin{tikzcd}
1 \arrow[r] & \Gamma^1 \arrow[r] \arrow[d, hook] \arrow[dr] & \Gamma^0 \arrow[r] & N(H)/H \\
 & H & \mathrm{Aut}_H(\pi_e) &
\end{tikzcd}

\note{la flèche verticale $\Gamma^1\to H$ est tracée barrée d'un petit trait ; on la rend par une inclusion, lecture incertaine.}
\drawing{un sommet $e$ d'où partent une arête vers un autre sommet et une flèche $f$ ; à côté, l'image de la situation en $\gamma e$.}

\struck{\ill{}}

$K=(K_0,K_1,\ K_1\to K_0\ \text{surj})$ (graphe \add{\uncertain{connexe}}), $G$ groupe, \struck{d'automorphismes de} \add{opérant sur} $K$, on suppose que $G$ est transitif sur l'ens\supplied{emble} $K_1$ des arêtes de $K$.
\[
\begin{matrix}
K_1\simeq G/H_1 \\
\downarrow \\
K_0=G/H_0
\end{matrix}
\qquad \sigma\in N(H)/H \qquad \sigma^2=1
\]
$C=(C_0,C_1\dots)$ \emph{groupoïde}, $G$ opère sur $C_0$
$K\hookrightarrow C$ $G$-morphisme ou \uncertain{2-morph\supplied{isme}}, $K_0\xrightarrow[\sim]{\text{bij}} C_0$
\drawing{à gauche, trois petites figures : une arête orientée de $e$ vers $e'$ ; une arête issue de $\gamma e$ ; une boucle en un sommet, avec une flèche sortante.}
\[
\begin{matrix}
1\to\Gamma^0\to\Gamma\to\mathrm{Aut}(\mathcal{C})\\
*\\
\Gamma\to G \qquad 1\to\Gamma^1\to\Gamma^0\to N(H)/H_0\\
*\\
1\to\Gamma^2\to\Gamma^1\hookrightarrow\mathrm{Aut}_G\,\pi_e\\
*\\
\pi_e,\ \mathrm{Hom}(e,e')\qquad \pi_e\\
\shortparallel\\
T_{e,e'} \qquad N_{F(\pi_e)}H_0
\end{matrix}
\]
\note{dans la deuxième suite, un $G$ entre $\Gamma^0$ et $N(H)/H_0$ est biffé ; la troisième suite (en $\Gamma^2$) est barrée de trois traits obliques.}

\note{le bas de la page est un brouillon de suites exactes superposées ; la disposition ci-dessus suit l'ordre des lignes, sans prétendre restituer l'alignement des étoiles, placées par lui sous $\mathrm{Aut}(\mathcal{C})$. « $N_{F(\pi_e)}H_0$ » est une lecture incertaine d'un groupe surchargé.}

\page{59}
\drawing{page de figures seules, reprenant celle de la p.~52 : en haut, un digone hachuré (deux arêtes joignant deux sommets) et un triangle à trois sommets marqués ; à droite, le sommet $s_0$ avec les arêtes $a_0$ (vers $s_1$), $a_1$ (vers $s_2$), $a_2$ (vers $s_3$) et les faces $f_0$, $f_1$ ; au centre, la figure développée : de $s_0$ partent $a_0$ vers $s_1$, $a_1$ vers $s_2$, $a_2$ vers $s_3$ et une quatrième arête vers le bas ; faces $f_0$ (hachurée, entre $a_0$ et $a_1$), $f_1$ (entre $a_2$ et $a_0$), $f_2$ (sous $a_1$), $f_3$ au-delà de $s_1$, et en $s_1$ l'arête $a_2'$. Le reste de la page est blanc.}

\end{document}
